% Chapter 49: Relativistic Quantum Mechanics
\chapter{Relativistic Quantum Mechanics}

Previous chapters gave us two formidable weapons. Special relativity binds energy, momentum, and mass through \(E^2=p^2c^2+m^2c^4\) in Chapter 40. Quantum mechanics describes microscopic states by wavefunctions evolving under Schrödinger's equation in Chapters 43 and 44. Each has been unfailingly successful on its own battlefield. What happens if we stitch them together? Sparks fly from the seam, producing three unrequested things: antimatter, spin, and a worldview that rewrites what a particle means.

\step{A Counterintuitive Opening}

A hospital PET scan, positron emission tomography, commonly images tumors and brain function. Its routine sounds ordinary: a nurse injects a fluorine-containing glucose analog into your vein, you lie quietly in a ring-shaped machine, and ten or so minutes later receive a color image of metabolic activity throughout your body.

That injection is anything but ordinary. Its fluorine-18 isotope spontaneously emits a positron, exactly an electron's mass but opposite charge, dramatically named the electron's antimatter. Traveling less than a millimeter through tissue, it meets an ordinary electron. Then something defies childhood intuition: both disappear without debris, leaving two gamma rays traveling almost exactly opposite ways, each carrying 511 keV. A detector ring records their precise arrival times, and a computer reconstructs millions of origins to draw the image.

Behind your medical report, antimatter forms inside you and annihilates with matter into light. School says matter cannot simply disappear and chemical reactions conserve atoms, yet two real particles vanish, leaving light. Stranger still, this particle was predicted before discovery: in 1928 a twenty-seven-year-old theorist trying to combine quantum mechanics and relativity found an extra solution in his equation. Theory declared it should exist; four years later an experimentalist found it in cloud-chamber photographs.

How can an equation predict an unseen particle? Why does joining two flawless theories first produce contradictions before revealing a new continent? That is our route.

\step{Make a Guess First}

Record your predictions. Faced with two successful theories, physicists share your instinct: surely a patch will do.

First, Schrödinger's kinetic term \(p^2/2m\) is only a low-speed approximation. Replace it with the correct relativistic version: does that produce a satisfactory relativistic Schrödinger equation? Choose (a) straightforward success, (b) minor repairable problems, or (c) an unpatchable fundamental wall.

Second, can a near-light-speed electron still be treated as one particle, with one wavefunction and a tracked probability distribution as in Chapter 43? When collisions become energetic enough, is the world's particle count a fixed roster or one subject to revision?

Third, if a new theory predicts an unheard-of particle, do physicists immediately trust the equation or first suspect a mistake? Think of finding a negative solution when a problem asks how many people there are.

Write your answers in the margin. By the end, the first answer approaches (c), the second is a revisable roster, and even Dirac hesitated three years over the third.

\step{Hands-on Experiment}

These phenomena occur at the 511 keV scale, so particle annihilation cannot be demonstrated directly at home. But a basin of water echoes the central conceptual reversal: particle number is not conserved; field excitations are fundamental.

\begin{experiment}[Counting Particles in a Basin]
Materials: a large flat-bottomed dish or sink, clean water, a chopstick, and a coin.

Step one: let the surface settle completely. Tap it quickly with the chopstick tip. A neat ripple ring travels outward, a wave packet resembling a little object hurrying along.

Step two: slap the surface with your whole palm. Now many ripples of different sizes and directions pass through one another, split, and merge. A stronger slap produces more numerous, varied patterns.

Step three: put a coin on the bottom to create a shallow patch and tap again. Ripples diffract and split above the coin, then superpose beyond it.

Count the waves. There is no definite answer: one packet can disperse into many ripples, and many can concentrate into one bulge. Wave number is not conserved; the conserved account tracks where your input energy went.
\end{experiment}

The analogy resembles our eventual worldview: water is the field, ripples its excitations, their number changing with supplied energy, without a law conserving total ripples. It differs because water-wave energy is continuously divisible, whereas quantum-field excitations come in packets whose size is fixed by Planck's constant. Water waves decay into heat; quantum excitations can genuinely appear and disappear in pairs out of nothing. This basin lends intuition, not evidence.

% BEGIN TRANSLATOR SAFETY NOTE
\textbf{Translator safety note:} Do not improvise the cloud-chamber project below. Use a supervised educational procedure: dry ice needs cold-rated gloves, goggles, ventilation, and gas-venting storage, never a sealed dry-ice container. Alcohol is flammable; keep it away from flames, sparks, and heat. See \href{https://stacks.cdc.gov/view/cdc/103603/cdc_103603_DS1.pdf}{CDC dry-ice guidance} and \href{https://www.cdc.gov/niosh/npg/npgd0359.html}{NIOSH isopropyl-alcohol hazards}.
% END TRANSLATOR SAFETY NOTE

If circumstances permit, a closer home project makes a simple Wilson cloud chamber with dry ice and alcohol. Numerous online tutorials exist; guard against dry-ice frostbite. Shine a flashlight into the darkened chamber and watch occasional white tracks from cosmic rays, mostly \(\mu\) particles. Formed high in the atmosphere, their rest lifetime is only about 2.2 microseconds; relativistic time dilation lets them reach the ground, as Chapter 38 explained. Here the problem is not how long they survive but their decay into other particles: kinds and counts change before your eyes. A theory admitting only a fixed particle count cannot even copy the cast list.

\step{The Old Model Runs into Trouble}

Like physicists in the 1920s, replace Schrödinger-equation parts carefully and see what jams.

First, time and space have unequal status. Chapter 44's equation contains a first time derivative but a second spatial derivative, the wavefunction's spatial curvature. Time is privileged. Yet Chapter 39 requires relativity to mix time and space between reference frames. An equation treating them differently changes form under moving rulers. It resembles a map using meters north--south and feet east--west: one user manages, but comparing maps demands peculiar conversions everywhere. This is structural failure, not cumbersome arithmetic.

Second, replacement produces negative probabilities. Translate \(E^2=p^2c^2+m^2c^4\) into operators, as the mathematical bridge explains. The resulting Klein--Gordon equation is second order in both time and space, curing Lorentz symmetry but introducing two illnesses. Taking square roots gives branches \(E=+\sqrt{p^2c^2+m^2c^4}\) and \(E=-\sqrt{p^2c^2+m^2c^4}\). What is energy below nothing? Why would particles not fall toward negative infinity, releasing unlimited energy? Worse, the proposed probability of finding a particle somewhere can become negative. Negative probability is meaningless, like drawing minus three winning tickets. This is no arithmetic mistake but a flaw in the translation. The equation later proved valuable for spin-zero particles, assigned not to a single-particle wavefunction but a quantum field. Sometimes an equation fails only because it has the wrong job.

Third and fundamentally, both repairs silently assume we always describe the same particle. But relativity gives \(E=mc^2\): energy and mass are exchangeable accounts. Once collision energy exceeds \(2mc^2\), two colliding particles can create an additional pair. Accelerators do it daily; even two energetic photons can create an electron--positron pair in the Breit--Wheeler process, directly observed by STAR in 2021 using intense photon fields around near-light-speed gold nuclei. Conversely, a pair can annihilate into light. A theory insisting there is forever one particle cannot write even the first page of this ever-changing cast's script.

Together these difficulties demand rebuilding foundations, not renovating an old equation by swapping parts.

\step{Inventing New Concepts}

In 1928 Paul Dirac reversed the approach: retain Schrödinger's first-order time structure, necessary for the present state to determine the next while preserving nonnegative conserved probability, but require relativity too. What equation could satisfy both?

The conditions looked impossible. Squaring a first-order time equation must recover \(E^2=p^2c^2+m^2c^4\). Ordinary coefficients cannot make \((ap+b)^2\) equal \(p^2c^2+m^2c^4\) without unwanted cross-terms. Dirac's startling move replaced numbers with matrices. The wavefunction became four numbers in a column, a spinor; coefficients became four \(4\times4\) matrices obeying special multiplication rules. This was not virtuosity for its own sake but a structure forced between two principles. The mathematical bridge will expose its compact logic.

Then the equation began speaking for itself, delivering three unexpected gifts far beyond Dirac's order.

First, spin appeared automatically. Solutions intrinsically carry angular momentum \(\hbar/2\), without Chapter 47's extra experimental postulate; the electron's magnetic moment, with \(g\) factor about 2, emerges too. The six-year-old Stern--Gerlach puzzle of two silver beams had been hiding in matrix structure. A theory invented for another purpose solves an old case free of charge, a strong historical sign of being on track.

Second, negative energies remained, but Dirac kept them. His audacious response was a historical interpretation, not an experimental fact: all negative-energy states are already filled by electrons in an invisible Dirac sea. Pauli exclusion prevents ordinary electrons from falling in, stabilizing the world. Knocking out one electron leaves a hole behaving exactly like an equal-mass positively charged particle. Intuition stumbled: Dirac initially identified it with the proton, then the only known positive particle. But protons weigh 1836 electrons, incompatible with the required equal mass. Theory had to predict a new particle. In 1932 Anderson photographed it in a cosmic-ray cloud chamber: curvature indicated positive charge and electron-like lightness. The positron was found. The sea is now dispensable historical scaffolding; the positron remains established experimental fact.

Third, conservation laws were reshuffled. Particle-number conservation retired: particles and antiparticles can appear or annihilate in pairs. Deeper conservation of energy, momentum, and charge remained clearer than before. A hole's opposite charge ensures creating or annihilating a pair preserves zero total charge. The conserved quantity is not how many things exist but a more abstract account.

The new concept is emerging: fully reconciling quantum mechanics with relativity costs us the belief that particles themselves are the most fundamental objects.
